% !PRINTED 350
\weberpage{350}
%<<CONT
be conjugate pairs. But then the three roots $(1,1)$ $(2,0)$ $(0,2)$ form a triple, whose conjugates $(1,1)$ $(1,0)$ $(0,1)$ form no triple, which is impossible. Each of the three real roots $(0,0)$, $(1,1)$, $(2,2)$ therefore occurs, besides in (1), in one and only one real triple, whose two other roots are conjugate imaginary; and the third case thus leads to the arrangement of the real and imaginary roots that we have come to know in the points of inflection of the curve of the third order.

But that the cases 1. and 2. can also occur is shown by the following consideration: From a general equation of degree $9$ one obtains a triple equation by adjunction of a function belonging to the linear group $SQ$. Such a function is, e.g.:

\[
\begin{aligned}
v &= (0,0)(1,1)(2,2) + (0,0)(1,2)(2,1) + (0,0)(1,0)(2,0) \\
&\quad + (0,0)(0,1)(0,2) + (1,1)(1,2)(1,0) + (1,1)(2,1)(0,1) \\
&\quad + (1,1)(0,2)(2,0) + (1,2)(0,1)(2,0) + (2,2)(2,1)(2,0) \\
&\quad + (2,2)(1,2)(0,2) + (2,2)(0,1)(1,0) + (2,1)(0,2)(1,0),
\end{aligned}
\]

and this function is real in the cases 1., 2., 3. Hence the extended domain of rationality, in which our equation is a triple equation, is also real.

\begin{center}
\rule{0.22\textwidth}{0.4pt}
\end{center}
